\documentclass[12pt]{article}
\usepackage{text_style}
\usepackage{a4wide}
\title{Diffusion Processes for Inference of Structural Causal Models}
\author{IA pour la Science}
\date{7 May 2026}
\begin{document}
\maketitle

\section{Linear models excluding State-Space models}

List of SCM-model parameters:
\begin{enumerate}
\item[] $\bw$ the linear parameters of the model,
\item[] $\bv$ the non-linear parameters of the model,
\item[] the covariance (mutual information) of the variables $X_1,\dots,X_n$,
\item the covariance of the linear part of the parameters $\bw$,
\item the covariance of the non-linear part of the parameters $\bv$.
\end{enumerate}

The transformations
\begin{enumerate}
    \item only G to f it admissible (there is no conversion from the model to G)
\end{enumerate}

Architectures of the linear models:
\begin{enumerate}
    \item linear model $\hat{\bx}=\bW\T\bx$: a single matrix (2-index) that maps $X_1,\dots,X_n$ to $\hat{X}_1,\dots,\hat{X}_n$,
    \item generalized linear model $\bsig\circ\bW\T\bx$: non-parametric link function of the linear model,
    \item special non-linear model $g_1(\bv_i) \circ g_r(\bv_r)(\bx)$: a model with functions $g_i(\bv_i)(X_j)$,
    \item general non-linear model $\bW\T g_i(\bv_i)(X_j)$ (obsoleted): a linear combination of the  model with functions,
    \item generalized non-linear model  $\bsig \circ \bW\T g_i(\bv_i)(X_j)$: a non-linear non-parametric transformation of the linear combination
\end{enumerate}

Optimization:
\begin{enumerate}
\item for linear models is MLS,
\item for non-linear models is Levenberg-Marquardt,
\item for the non-linear models is Levenberg-Marquardt,
\item for the rest it will 
\end{enumerate}


Here the notation $\circ$ means the composition of functions.
It could be replaced by brackets.



%\end{document}
%=======================================================
\section{Ideal Magnetohydrodynamics (MHD) Equations}

The equations

\begin{align}
\frac{\partial \rho}{\partial t}
+ \nabla \cdot (\rho \mathbf{v})
&= 0,
\\[1ex]
\frac{\partial (\rho \mathbf{v})}{\partial t}
+ \nabla \cdot
\left(
\rho \mathbf{v}\mathbf{v}
-
\mathbf{B}\mathbf{B}
\right)
+ \nabla p
&= 0,
\\[1ex]
\frac{\partial \mathbf{B}}{\partial t}
-
\nabla \times (\mathbf{v}\times\mathbf{B})
&= 0,
\end{align}

are a simplified form of the \emph{ideal magnetohydrodynamics (MHD)} equations.
They describe the dynamics of an electrically conducting fluid
(plasma, ionized gas, liquid metal, etc.) interacting with a magnetic field.

The primary variables are

\begin{itemize}
\item $\rho(\mathbf{x},t)$: mass density,
\item $\mathbf{v}(\mathbf{x},t)$: fluid velocity,
\item $\mathbf{B}(\mathbf{x},t)$: magnetic field,
\item $p(\mathbf{x},t)$: pressure.
\end{itemize}

\subsection{Continuity Equation}

\begin{equation}
\frac{\partial \rho}{\partial t}
+
\nabla\cdot(\rho\mathbf{v})
=
0.
\end{equation}

This equation expresses conservation of mass.

The first term,

\[
\frac{\partial \rho}{\partial t},
\]

measures the rate of change of density at a fixed point in space.

The second term,

\[
\nabla\cdot(\rho\mathbf{v}),
\]

represents the divergence of the mass flux.
It measures the net amount of mass flowing out of a small control volume.

If more mass leaves a region than enters it, the density decreases.
If more mass enters than leaves, the density increases.

Thus the continuity equation states that mass can neither be created nor destroyed.

\subsection{Momentum Equation}

\begin{equation}
\frac{\partial (\rho\mathbf{v})}{\partial t}
+
\nabla\cdot
\left(
\rho\mathbf{v}\mathbf{v}
-
\mathbf{B}\mathbf{B}
\right)
+
\nabla p
=
0.
\end{equation}

This equation is the fluid analogue of Newton's second law.

The quantity

\[
\rho\mathbf{v}
\]

is the momentum density of the fluid.

The term

\[
\frac{\partial(\rho\mathbf{v})}{\partial t}
\]

represents the local time rate of change of momentum.

The tensor

\[
\rho\mathbf{v}\mathbf{v}
\]

describes transport of momentum by fluid motion.
It accounts for momentum carried from one region to another.

The tensor

\[
\mathbf{B}\mathbf{B}
\]

represents magnetic stresses.
Magnetic field lines behave somewhat like stretched elastic strings:
they exert tension along their length and pressure perpendicular to themselves.

The pressure-gradient term

\[
\nabla p
\]

describes the force exerted by pressure differences.
Fluid tends to move from regions of higher pressure toward regions of lower pressure.

Overall, the momentum equation states that the fluid momentum changes because of
pressure forces and magnetic forces.

\subsection{Induction Equation}

\begin{equation}
\frac{\partial \mathbf{B}}{\partial t}
-
\nabla\times(\mathbf{v}\times\mathbf{B})
=
0.
\end{equation}

This equation governs the evolution of the magnetic field.

The quantity

\[
\mathbf{v}\times\mathbf{B}
\]

describes the interaction between fluid motion and the magnetic field.
Its curl determines how the magnetic field changes in time.

A fundamental consequence of this equation in ideal MHD is the
\emph{frozen-in flux property}:

\begin{quote}
Magnetic field lines move with the fluid.
\end{quote}

If the plasma stretches, twists, or compresses,
the magnetic field lines stretch, twist, or compress with it.

\subsection{Additional Equations}

A complete ideal MHD system normally includes the constraint

\begin{equation}
\nabla\cdot\mathbf{B}=0,
\end{equation}

which expresses the absence of magnetic monopoles.

An equation of state or energy equation is also required to determine the pressure.
For example,

\begin{equation}
p=(\gamma-1)\rho e,
\end{equation}

where $e$ is the specific internal energy and $\gamma$ is the adiabatic index.

Without such a relation the system is not closed.

\section{Where is Space in the Equations?}

The spatial coordinates are implicit in the notation.

More explicitly, each field depends on both position and time:

\begin{align}
\rho &= \rho(\mathbf{x},t),
\\
\mathbf{v} &= \mathbf{v}(\mathbf{x},t),
\\
\mathbf{B} &= \mathbf{B}(\mathbf{x},t),
\\
p &= p(\mathbf{x},t),
\end{align}

where

\begin{equation}
\mathbf{x}=(x,y,z).
\end{equation}

For example, the continuity equation can be written explicitly as

\begin{equation}
\frac{\partial \rho(x,y,z,t)}{\partial t}
+
\nabla\cdot
\left(
\rho(x,y,z,t)\,
\mathbf{v}(x,y,z,t)
\right)
=
0.
\end{equation}

The dependence on $(x,y,z,t)$ is usually omitted for brevity.

\subsection{Gradient Operator}

The gradient of a scalar field $p$ is

\begin{equation}
\nabla p
=
\left(
\frac{\partial p}{\partial x},
\frac{\partial p}{\partial y},
\frac{\partial p}{\partial z}
\right).
\end{equation}

It measures how rapidly the pressure changes from one point to another.

\subsection{Divergence Operator}

For a vector field

\[
\mathbf{F}
=
(F_x,F_y,F_z),
\]

the divergence is

\begin{equation}
\nabla\cdot\mathbf{F}
=
\frac{\partial F_x}{\partial x}
+
\frac{\partial F_y}{\partial y}
+
\frac{\partial F_z}{\partial z}.
\end{equation}

It measures the net outflow from an infinitesimal volume.

For the mass flux,

\begin{equation}
\nabla\cdot(\rho\mathbf{v})
=
\frac{\partial(\rho v_x)}{\partial x}
+
\frac{\partial(\rho v_y)}{\partial y}
+
\frac{\partial(\rho v_z)}{\partial z}.
\end{equation}

\subsection{Curl Operator}

The curl of a vector field is

\begin{equation}
\nabla\times\mathbf{F}
=
\begin{pmatrix}
\displaystyle
\frac{\partial F_z}{\partial y}
-
\frac{\partial F_y}{\partial z}
\\[2ex]
\displaystyle
\frac{\partial F_x}{\partial z}
-
\frac{\partial F_z}{\partial x}
\\[2ex]
\displaystyle
\frac{\partial F_y}{\partial x}
-
\frac{\partial F_x}{\partial y}
\end{pmatrix}.
\end{equation}

It measures the local rotational tendency of the field.

\subsection{One-Dimensional Example}

If the system varies only in the $x$ direction, then

\[
\rho=\rho(x,t),
\qquad
v=v(x,t).
\]

The continuity equation reduces to

\begin{equation}
\frac{\partial \rho}{\partial t}
+
\frac{\partial(\rho v)}{\partial x}
=
0.
\end{equation}

In this form the independent variables are explicit:

\begin{itemize}
\item $x$ is the spatial coordinate,
\item $t$ is time.
\end{itemize}

The full MHD equations generalize this idea to three spatial dimensions by using
the compact differential operators $\nabla$, $\nabla\cdot$, and $\nabla\times$.

Therefore, space enters the equations through the position vector

\[
\mathbf{x}=(x,y,z),
\]

and through the spatial derivatives acting on the fields.
The variables $\rho$, $\mathbf{v}$, $\mathbf{B}$, and $p$
are defined at every point in space and evolve in time according to the MHD equations.

%======================================








Two aspects of the SCM inference were discussed the last week:
\begin{enumerate}
\item The inference avoids the  Integer Liner Programming problem by replacing the probabilities of edges
\[\bGam \in [0,1]^{n\times n} \mapsto \bZ \in \{0,1\}^{n\times n}.\]
for the adjacency of edges. \emph{NB: it is solved below.}
\item  The list of variables is not fixed.
For the set of observable variables $X_1,\dots,X_n$
the model has to forecast the list of variables $X_1,\dots,X_n, \dots, X_{n^\prime}$.
By default, the value $n^\prime=0$. \emph{NB: it is under consideration.}
\end{enumerate}

The foundation model includes the elements of SCM selection:
\begin{enumerate}[1)]
\item maximizing likelihood of the DAG at each step,
\item do-intervention actions each step on the sampled node.
\end{enumerate}

\subsection*{The foundation model for the SCM inference}
The main type of the foundation model carries the (Riemannian) diffusion process.

The diffusion process involves the step~$t$ and two distributions to approximate the distribution of 
the adjacency matrix~$\bZ$.
Use the vector form~$\bz = \text{vec}(\bZ)$.
Denote the training distribution~$q(\bz)$ and the sampling distribution~$p(\bz)$.
During the model training one has to minimize the distance between two distributions at each step. 

The diffusion process invokes the local structural modeling:
\[
\bZ \mapsto f_g \mapsto \eta \mapsto \bw,
\]
with $\eta$ is the optimization algorithm that delivers the optimal parameters to the SCM at each step of the diffusion process.

The \emph{fine-tuning} algorithm~$\eta$ optimizes the parameters according to the diagram

% Optimization and inference by an SCM.
%{\color{red}\footnote*{\color{gray}These notations are equivalent: $x_i, x(i), i\to x$.}}
\begin{equation*}
\xymatrix{
\eta(-\lambda^2\nabla_{\textbf{w}}\loss+\textbf{w}^\prime)\ar[r]^{\qquad \eta} & (\hat{\bw}, {\bx}) \ar[r]^{f_G} &  \hat{\bx}\ar[r] & \loss(\mathbf{\textbf{w}\mid\bx,\hat{\bx}}) \ar@/_3pc/[lll]\\
%1 & 2 &  & y\ar[ur] & 5 \\
(t,s)\ar[r] & i\ar[r] & \bx\ar[ul]\ar[ur] & }
\end{equation*}
This diagram reads: the optimization algorithm $\eta$ estimates the SCM parameters~$\bw$ by the gradient of its loss and the previous parameters to solve the problem 
\begin{equation*}
\hat{\textbf{w}}= \arg\min \loss(\bw) \quad\text{given}\quad f_G(\hat{\bw},\bx)\mapsto \hat{\bx}.
\end{equation*}  
The algorithm $\eta$ uses some initial value~$\bw_0$ and a stop-criterion for its iterations.

\subsection*{A basic diffusion process for the data}
The data~$\bx$ only, without the binary adjacent $\bz$ for the model $f_G$ are sampled and inferred with the generative process below.

\begin{tikzpicture}[
        node distance=1.5cm,
    every node/.style=,
    var/.style={circle, draw, minimum size=1.2cm},
    arrow/.style={->, thick},
    curved/.style={->, thick, bend left=30},
]
%    node distance=1.5cm,
%   every node/.style={circle, draw, minimum size=1.2cm},
%    arrow/.style={->, thin},
%    dashedarrow/.style={->, thin, dashed, red}


% Nodes
\node (xT) {$\bx_T$};
\node (dots1) [right=of xT, draw=none] {$\cdots$};
\node (xt) [right=of dots1] {$\bx_t$};
\node (xtm1) [right=2.5cm of xt] {$\bx_{t-1}$};
\node (dots2) [right=of xtm1, draw=none] {$\cdots$};
\node (x0) [right=of dots2] {$\bx_0$};

% Forward arrows (top)
\draw[arrow] (xT) -- (dots1);
\draw[arrow] (dots1) -- (xt);
\draw[arrow] (xt) to node[below]{$p_\bth(\bx_{t-1}\mid \bx_t)$}(xtm1);
\draw[curved, bend right=60] (xt) to node[below] {$q(\bx_{t-1} \mid \bx_t)$} (xtm1);
\draw[curved, bend right=60] (xtm1) to node[above] {$q(\bx_{t} \mid \bx_{t-1})$} (xt);
\draw[arrow] (xtm1) -- (dots2);
\draw[arrow] (dots2) -- (x0);

\end{tikzpicture}

Given the data distribution $\bx_0\sim q(\bx)$ assert:
\begin{enumerate}[1)]
\item forward diffusion process $\bx_t = \sqrt{1-\beta_t} \bx_{t-1} + \sqrt{\beta_t} \bx_t$, sampling i.i.d. 
\[\bx_1,\dots, \bx_T\sim\cN (\boldsymbol{0}, \mathbf{I}),\]
\item sampling $q(\bx_{t-1}\mid\bx_t)$ and learning the parameters $\bth$ of $f_\text{ai}$: 
$p_\bth(\bx_{t-1}\mid\bx_t)$,  
\item reverse diffusion process 
\[
p_\bth(\bx_{0:T}) = p(\bx_T) \prod^T_{t=1} p_\bth(\bx_{t-1} \mid \bx_t),
\]
\item  slow learning gives
\[
p_\bth(\bx_{t-1} \mid \bx_t) = \cN(\bx_{t-1}; \bmu_\bth\bigl(\bx_t, t), \boldsymbol{\Sigma}_\bth(\bx_t, t)\bigr).
\]
\end{enumerate}

\subsection*{Coherent generation of SCM submodels with do-intervention}
%\subsection*{The do-intervention and SCM sub-modeling}
Each step of the diffusion delivers a set of submodels to reconstruct. The dotted lines of the diagram are sampling; the rest are inference or optimization.

\bigskip
\begin{tikzpicture}[
    node distance=2.5cm,
    every node/.style={font=\large},
    var/.style={circle, draw, minimum size=1.2cm},
    arrow/.style={->, thick},
    curved/.style={->, thick, bend left=30},
]

% Top latent chain
\node[var] (zt) {$\bz_t$};
\node[var, right=of zt] (ztm1) {$\bz_{t-1}$};

\node[left= 1cm of zt] (ztTm1) {$\dots$};
\node[var, left= 2cm of zt] (ztT) {$\bz_T$};
\node[right= 1cm of ztm1] (ztm0) {$\dots$};
\node[var, right= 2cm of ztm1] (zt0) {$\bz_0$};
\draw[arrow] (ztTm1) -- (zt);
\draw[arrow] (ztm1) -- (ztm0);

% Curved bidirectional arrows (approximation of sketch)
\draw[arrow] (zt) to node[above]{$p_\bth(\bz_{t-1}\mid \bz_t)$}(ztm1);
\draw[dotted, curved, bend right=40] (ztm1) to node[above] 
	{$q(\bz_t\mid \bz_{t-1})\underbrace{q(\bz_{t-1}\mid \bA_{t-1})}_{\text{MC/M-H prior}}$}(zt);
\draw[dotted, curved, bend right=40] (zt) to node[below] {$q(\bz_{t-1}\mid \bz_t)$} (ztm1);

% Middle variables
\node[var, below=2.5cm of zt] (wt) {$\bw_t$};
\node[var, below=2.4cm of ztm1] (wtm1) {$\bw_{t-1}$};
  
\draw[arrow] (wtm1) --  node[above]{{\small fine tune} $\loss$} (wt);
% Reconstruct the model parameters
\draw[dotted, curved, bend left=40] (wt) to node[above] {$q_\bw(\bw_{t-1}\mid \bw_t)$} (wtm1); 

% Vertical connections (adjoint operators)
\draw[arrow] (zt) -- node[left] {$f_G(t)$} (wt);
\draw[arrow] (ztm1) -- node[right] {$f_G(t-1)$} (wtm1);

% Bottom variables
\node[var, below=2.5cm of wt] (xt) {$\bx_t$};
\node[var, below=2.35cm of wtm1] (xtm1) {$\bx_{t-1}$};

\draw[arrow] (xtm1) -- node[above] {$\beps$} (xt);

% Vertical arrows (forward operators)
\draw[arrow] (xt) -- node[left] {$\eta$} (wt);
\draw[dotted, arrow, bend right=40] (xt) to node[left] {\raisebox{-12ex}{$\bA_t^\prime$}} (zt);
\draw[arrow] (xtm1) -- node[right] {$\eta$} (wtm1);
\draw[dotted, arrow, bend right=40] (xtm1) to node[right] {\raisebox{-12ex}{$\bA_{t-1}^\prime$}} (ztm1);

% Data terms at bottom
\node[below=1.2cm of xt] (dxt) {$\text{do}(\bx_t)$};
\node[below=1.2cm of xtm1] (dxtm1) {$\text{do}(\bx_{t-1})$};

\draw[arrow] (dxt) -- (xt);
\draw[arrow, bend right=40] (wt) to node[left] {\raisebox{4.5ex}{$\bA_t$}} (dxt);
\draw[arrow] (dxtm1) -- (xtm1);
\draw[arrow, bend right=40] (wtm1) to node[right] {\raisebox{4.5ex}{$\bA_{t-1}$}} (dxtm1);

%%%%%%%%%%% Part extra: add T and t=0 to the left and right of the scheme
% The second row, w
\node[left= 1cm of wt] (wtTm1) {$\dots$};
\node[var, left= 2cm of wt] (wtT) {$\bw_T$};
\draw[arrow] (wt) -- (wtTm1);
% right part
\node[right= 1cm of wtm1] (wtm0) {$\dots$};
\node[var, right= 2cm of wtm1] (wt0) {$\bw_0$};
\draw[arrow] (wtm0) --  (wtm1);
% The third row, x
\node[left= 1cm of xt] (xtTm1) {$\dots$};
\node[var, left= 2cm of xt] (xtT) {$\bx_T$};
\draw[arrow] (xt) -- (xtTm1);
% right part
\node[right= 1cm of xtm1] (xtm0) {$\dots$};
\node[var, right= 2cm of xtm1] (xt0) {$\bx_0$};
\draw[arrow] (xtm0) -- (xtm1);

\end{tikzpicture}
%\end{document}
\paragraph{The algorithm of SCM inference} includes the U-Net $f_{\text{ai}}(\bth)(\bx,\bz)$ to learn its parameters~$\bth$ and make an inference of the SCM $f_G$ along with its parameters $\bw$. The estimated parameters~$\hat{\bw}$ are inherited at each step~$t$. So that only fine-tuning optimization~$\eta$ goes to minimize $\loss(\bw_{t-1})\mapsto\bw_t$. %The data collections $\{\bX\}$ (denoted here by one sample $\bx$) couples the stricture $\bz$ with a transformer $(\bX,\bz)$. 
(The dataset refers to the stricture  $(\bX,\bz)$ with a transformer.) 

\begin{enumerate}
\item Sample a new SCM adjacency $\bz_t$.
\item Construct the new model $f_G$ and its initial parameters $\bw_t$ from the step $t-1$.
\item Fine-tune the SCM parameters $\hat{\bw}_t = \arg\min \loss(\bw_t \mid \hat{\bw}_{t-1})$ with the initial parameters.
\item Estimate the covairiance $\hat{\bA}^{-1} = \textbf{cov}(\bw)$ using the hessian of $\loss$, the Fisher matrix.
\item Apply the CI-criterion $\rho(\bA_t)$  to make the intervention $\doc(\bx_t)$.
\item Make the intervention, sampling $\doc(\bx_t)$ with the intervention parameters $\bA^\prime_t$
\item Compute the partial residuals $ \bx_{t-1} - \hat{\bx}_{t-1}= \beps_t \mapsto \bx_t$.
\item Learn the $f_\text{ai}$ parameters by reconstructing $p_\bth(\bz_{t-1} \mid \bz_t)$.
\item Go to the step 1 and sample the next adjacency $\bz_{t-1}$.
\end{enumerate}

\vspace{5cm}
\section*{List of the modules}
%remove the document; it is OLD!!!!
%\newpage
\subsection*{The empirical distribution of the observations}
$P(X_1,\dots,X_n)$ is the distribution of the observable data.
Its approximation is given by the Gaussian Process Regression model
\[
p(\bx \mid f_\text{gp})
\]
on the dataset~
\[\{\bx_i\in\mathbb{R}^n\}.\]
It is indexed by some function of the time and space $(t,s)\mapsto i$.

\subsection*{A model from the family of universal approximators}
For a sample $\bx$ and its reconstruction $\hat\bx$ the likelihood function provided a universal model~$f_\text{nn}(\bth,\bx)$ is
\[p(\bx\mid f_\text{nn}, \bth).\]
Without the a priori information about the model parameters~$\bth$ the corresponding loss function is
\[
\loss(\bth) = -\log p(\bx\mid f_\text{nn}, \bth) = (\hat{\bx} - \bx)\T \mathbf{B}^{-1}(\hat{\bx} - \bx).% + (\bth - \boldsymbol{\mu})\T \mathbf{A}^{-1}(\bth - \boldsymbol{\mu}),
    \]
The matrix $\mathbf{B}$ is the covariance of the reconstruction error.
This matrix in some statements depends on the index~$i$.
This model as a parametric one has a structure to make a distillation to the causal graph
\[
f_\text{nn}(\bth, \bx)  \overset{\text{extract}}{\longrightarrow} \bZ_\text{nn} \overset{\text{dist}}{\longrightarrow}  \bZ.
\]

\subsection*{An SCM we have to forecast}
A structural causal model
\[
    \hat{\bx} = f_G(\bw,\bx) \in\mathcal{F}_\text{scn}
\]
has a structure~$\bZ$ and its probabilistic representation $\bGam$.
The elements of the matrix~$\bGam$ are the probabilities of the causal relationships between the elements of the model structure~$f_G$.
The matrix~$\bZ$ is the adjacency matrix of the causal graph, the binary representation of~$\bGam$.

\medskip
Note: \emph{we suppose that SCM does not fit the whole $(s,t)$-data but a local part of it.
So we also call SCM a local model~$f_G$, in contrast to the universal model~$f_\text{nn}$.}

\subsection*{A foundational model to forecast the SCM}
There given pairs a sample $\bx_i$ and the corresponding SCM~$f_G$ with its structure $\{\bZ_i\}$.
Call the set of pairs
\[(\bx_i, \bZ_i), \quad i\in\cB
\]
a collection of the SCM forecasting.
This collection $\cB$ shows one case of mathematical modeling that results in an interpretable model~$f_G$.
This model fits samples $\bx_i$ from some local data~$i\in\cB$.

Call a model that forecasts an SCM $\mathcal{F}$ a \emph{foundation model}
\[
f_\text{ai}: \bx_\cB \mapsto f_G \in \mathcal{F}.
\]
The SCM is a parametric model
\[
    f_G = f_G (\bw,\bx) \quad\text{with the parameters}\quad \hat{\bw} = \arg\min_{\bw} \loss(\bw).
\]
The foundational model returns not the SCM~$f_G(\hat{\bw},\bx)$ itself,
but only a probabilities~$\bGam$ on the structure~$\bZ$.
So we rewrite the previous definition as
\[
f_\text{ai}: \bx_\cB \mapsto \bGam.
\]

\medskip
Note: \emph{Since a case of mathematical modeling, call it $\cB$ has a history of requests to the foundational model,
    the next information is available for it:
\begin{enumerate}[1)]
    \item the optimal parameters $\bw$ along with $\bZ, \bGam$ for the SCM $f_G$,
    \item the prior probabilities of the model parameters $p(f_G)$ or all models in the family $\mathcal{F}$,
    \item the other information about the case of mathematical modeling $\cB$ that is relevant for the forecasting of the SCM.
\end{enumerate}}

\subsection*{The Integer Linear Programming for the causal graph selection}
The matrix $\bGam$ comes along with the descriptions of its nodes $X_1,\dots,X_n$, and the probabilities.
The SCM $f_G$ is defined by the graph $G=(\bZ,X)$ by solving
\[
  \text{ILP}(\bGam,P_X) \to \bZ \to f_G.
\]

\subsection*{The evidence of the SCM}
A probabilistic description~$\bGam$ of a model $f_G$ from the class of \emph{admissible models} $\mathcal{F}$ is generated from a parametric distribution
\[
    \bGam \sim \text{Dir}(\bpi) \quad\text{generated by MCMC}\quad f_\text{mh}(\bpi),
\]
here the lowercase letters mean Metropolis-Hastings Algorithm for MCMC.
(A simple genetic algorithm may bring some decent results, too).
This distribution provides a probabilistic space for the similar models.
See the document on \emph{Multimodeling}.

The model $f_G$ and the corresponding graph has its likelihood
\[
    p(\bx|f_G).
\]
See the document on \emph{Causal graph likelihood}.

\medskip
Note: \emph{Two-level bayesian inference introduces the probability of the SCM parameters on the first level.
And the probability of the model on the second level is:
    \begin{enumerate}[1)]
    \item the prior of the SCM is delivered by the generative model $f_\text{mh}$
    \item the likelihood of the SCM is delivered by the forecasting model $f_\text{ai}$.
\end{enumerate}}

\subsection*{The criterion of causality}
The parameters of the model $f_G$ has their covariance matrix.
The prior of the model parameters is
\[
    - \ln p(\bw|f_G) = (\bw - \bmu)^{\mathsf{T}} \bA^{-1} (\bw - \bmu).
\]
The algorithm of the Bayesian hyperparameters $(\bmu, \bA)$ estimation
(and their connection between the causality)
is described in the document on \emph{parameter analysis}.

\subsection*{A model for the intervention}
See the document on the \emph{intervention strategy}.

\end{document}
